\chapter{Teoretyczne podstawy architektury mikroserwisowej}
\label{chap:podstawy-mikroserwisow}

Rozdział ten przedstawia teoretyczne podstawy architektury mikroserwisowej. Omówiono w~nim różnice
pomiędzy architekturą monolityczną a~mikroserwisową, kluczowe cechy mikroserwisów, modele komunikacji międzyusługowej, a~także wzorce projektowe
stosowane w~systemach rozproszonych. Rozdział zamyka omówienie konteneryzacji oraz roli orkiestracji,
stanowiące wprowadzenie do rozdziału~\ref{chap:kubernetes}.

\section{Architektura monolityczna a mikroserwisowa}
\label{sec:monolit-a-mikroserwisy}

Architektura monolityczna to model budowy aplikacji, w której całość logiki biznesowej kompiluje się i wdraża jako
pojedynczą jednostkę uruchomieniową. Model ten ma istotne zalety. Jest on prosty koncepcyjnie. Transakcje mogą być
lokalne i obejmować cały stan aplikacji. Uproszczone jest testowanie i debugowanie, a~wdrożenie sprowadza się do uruchomienia 
jednego artefaktu. Z~tych powodów monolit pozostaje racjonalnym wyborem dla małych zespołów i~systemów 
o~ograniczonej złożoności \cite{newman2021}.

Problemy architektury monolitycznej ujawniają się wraz z rozbudową systemu i~organizacji. Skalowanie możliwe
jest poprzez dodanie zasobów lub powielenie całej aplikacji. Nawet jeśli wąskim gardłem jest pojedynczy moduł.
Każda zmiana wymaga zbudowania i wdrożenia całej aplikacji. Co wydłuża cykl wydawniczy i~zwiększa ryzyko regresji.
Brak wymuszenia granic pomiędzy modułami może powodować z czasem erozje, prowadząc do silnego sprzężenia kodu.
Finalnie cały system związany jest jedną technologią i~jedną wersją zależności, a~awaria dowolnego fragmentu 
może unieruchomić aplikację w~całości \cite{richardson2018}.

Architektura mikroserwisowa stanowi odpowiedź na te ograniczenia. Lewis i~Fowler definiują ją jako podejście,
w~którym pojedyncza aplikacja budowana jest jako zbiór niewielkich usług. Każda działa we własnym procesie
i~komunikuje się z~pozostałymi za pomocą lekkich mechanizmów. Usługi zbudowane są wokół zdolności biznesowych
i~wdrażane niezależnie w~sposób w~pełni zautomatyzowany \cite{fowler2014}. Zestawienie najważniejszych różnic
pomiędzy oboma podejściami przedstawia tabela~\ref{tab:monolit-vs-mikroserwisy}.

\begin{table}[H]
    \centering
    \caption{Porównanie architektury monolitycznej i mikroserwisowej 
    (opracowanie własne na podstawie \cite{newman2021,richardson2018})}
    \label{tab:monolit-vs-mikroserwisy}
    \renewcommand{\arraystretch}{1.25}
    \setlength{\tabcolsep}{8pt}

    \begin{tabularx}{\textwidth}{>{\raggedright\arraybackslash\bfseries}p{3.4cm}
                                    >{\raggedright\arraybackslash}X
                                    >{\raggedright\arraybackslash}X}
        \toprule
        \textbf{Kryterium} 
        & \textbf{Architektura monolityczna} 
        & \textbf{Architektura mikroserwisowa} \\
        \midrule

        Jednostka wdrożenia 
        & Cała aplikacja wdrażana jako jedna całość 
        & Pojedyncza, niezależna usługa \\

        Skalowanie 
        & Skalowanie całej aplikacji, nawet gdy obciążony jest tylko jeden moduł 
        & Skalowanie niezależne, dostosowane do potrzeb konkretnej usługi \\

        Baza danych 
        & Najczęściej jedna wspólna baza danych 
        & Odrębna baza danych lub schemat danych dla każdej usługi \\

        Komunikacja wewnętrzna 
        & Bezpośrednie wywołania funkcji lub metod w ramach jednej aplikacji 
        & Komunikacja sieciowa, np. REST, komunikaty lub zdarzenia \\

        Spójność danych 
        & Zwykle zapewniana przez transakcje ACID 
        & Często oparta na spójności ostatecznej \\

        Dobór technologii 
        & Zazwyczaj jednolity stos technologiczny 
        & Możliwość niezależnego doboru technologii dla poszczególnych usług \\

        Izolacja awarii 
        & Ograniczona; awaria jednego modułu może wpływać na całą aplikację 
        & Większa; awaria może zostać ograniczona do pojedynczej usługi \\

        Złożoność operacyjna 
        & Relatywnie niska, prostsze wdrażanie i monitorowanie 
        & Wysoka, wymaga orkiestracji, monitoringu i zarządzania komunikacją \\

        Cykl wydawniczy 
        & Wspólny dla całej aplikacji, zwykle mniej elastyczny 
        & Niezależny dla poszczególnych usług, umożliwiający częstsze wdrożenia \\

        \bottomrule
    \end{tabularx}
\end{table}

Należy podkreślić, że mikroserwisy nie są rozwiązaniem uniwersalnie lepszym. Przenoszą one złożoność z~wnętrza aplikacji
do przestrzeni pomiędzy usługami. Takimi jak sieć, infrastruktura czy procesy wdrożeniowe. Wybór tej architektury jest
uzasadniony wówczas, gdy gdy korzyści z~niezależnego skalowania i~wdrażania przewyższają koszt zwiększonej
złożoności operacyjnej. Z tych powodów praktyczna realizacja mikroserwisów jest bardzo związana z~technologiami
chmurowymi i~ich automatyzacją.

\section{Charakterystyka mikroserwisów}
\label{sec:charakterystyka-mikroserwisow}

Literatura przedmiotu \cite{newman2021,richardson2018} wskazuje zestaw cech konstytuujących architekturę
mikroserwisową, które można sprowadzić do następujących zasad.

\textbf{Modelowanie wokół domeny biznesowej.} Granice usług powinny odpowiadać granicom kontekstów biznesowych
(\emph{bounded contexts}) w~rozumieniu projektowania sterowanego domeną (\emph{Domain-Driven Design}) \cite{evans2003}.
Usługa bierze odpowiedzialność za wyodrębniony fragment domeny. W realizowanym systemie te fragmenty to:
zamówienia, katalog produktów oraz płatności. Minimalizuje to liczbę zmian wymaganych do koordynacji między
usługami.

\textbf{Niezależność wdrożeniowa.} Każda usługa może być zbudowana, przetestowana i~wdrożona bez udziału pozostałych. 
Separacja ta pozwala rozwijać i testować aplikacje samodzielnie. Wymaga to jawnych wersjonowanych kontraktów komunikacyjnych 
oraz unikanie współdzielenia struktur danych. Bardzo dobrą praktyką jest rozdzielenie struktur na komunikacyjną i transportową.

\textbf{Ukrywanie informacji i~własność danych.} Usługa udostępnia na zewnątrz wyłącznie swój interfejs, ukrywając
szczegóły implementacji. Konsekwencją jest zasada \emph{database-per-service}. Każda usługa jest wyłącznym właścicielem swoich danych.
Inne usługi mogą uzyskać do nich dostęp jedynie poprzez jej interfejs. Złamanie tej zasady prowadzi do sprzężenia na poziomie schematu, 
które w~praktyce odbiera usługom niezależność wdrożeniową \cite{newman2021}.

\textbf{Decentralizacja i~autonomia technologiczna.} Usługi mogą być rozwijane przez niezależne zespoły 
i~realizowane w~różnych technologiach, dobranych do charakteru problemu.

\textbf{Odporność na awarie.} W systemie rozproszonym awarie są normalnym zjawiskiem. Architektura musi zakładać że usługi 
mogą być chwilowo niedostępne. Ograniczenie takich awarii można osiągnąć poprzez asynchroniczną komunikacje oraz samonaprawiające się usługi.

\textbf{Obserwowalność.} Zachowanie systemu rozproszonego wymaga dokładnego monitorowania, ponieważ każda usługa działa całkowicie oddzielnie.
Wymagane jest dokładne monitorowanie każdej z nich w celu wykrywania anomalii. Niezbędne więc staje się gromadzenie metryk,
logów i informacji o stanie każdej usługi. Zagadnienie to rozwinięto w~rozdziale~\ref{chap:kubernetes}
oraz, w~odniesieniu do implementacji, w~rozdziale~\ref{chap:implementacja-mikroserwisow}.

\section{Komunikacja synchroniczna i asynchroniczna}
\label{sec:komunikacja-synchroniczna-asynchroniczna}

Dekompozycja systemu na usługi czyni z~komunikacji międzyusługowej zagadnienie architektonicznie pierwszorzędne.
Wyróżnia się dwa zasadnicze modele tej komunikacji, różniące się sprzężeniem czasowym pomiędzy stronami.

\subsection{Komunikacja synchroniczna - REST}
\label{subsec:komunikacja-synchroniczna}

W~modelu synchronicznym usługa wywołująca wysyła żądanie i~wstrzymuje przetwarzanie do czasu otrzymania odpowiedzi.
Najpowszechniejszą realizacją tego modelu są interfejsy REST (\emph{Representational State Transfer}) oparte na
protokole HTTP i~formacie JSON. Zasoby identyfikowane są za pomocą adresów URL, a~operacje metodami HTTP (GET, POST, PUT, DELETE).
Zaletami podejścia jest prostota, łatwość testowania oraz zgodność z~modelem żądanie-odpowiedź.
Podobnym dla interakcji użytkownika z serwisem.

Wadą komunikacji synchronicznej jest sprzężenie czasowe. Obie strony muszą być dostępne jednocześnie, a~niedostępność
lub spowolnienie usługi wywoływanej propaguje się na wywołującą. W~łańcuchach wywołań prowadzi to do awarii
kaskadowych, w~których problem jednej usługi unieruchamia kolejne \cite{richardson2018}. Powoduje to złamanie paradygmatu
mikroserwisów czyli odporność na awarie. Z~tego względu komunikacje synchroniczną stosuje się głównie na styku użytkownik-system.

\subsection{Komunikacja asynchroniczna i architektura sterowana zdarzeniami}
\label{subsec:komunikacja-asynchroniczna}

W~modelu asynchronicznym nadawca przekazuje komunikat pośrednikowi brokerowi komunikatów i~natychmiast
kontynuuje przetwarzanie, nie oczekując na odbiorców. Szczególną postacią tego modelu jest architektura sterowana
zdarzeniami (\emph{event-driven architecture}), w~której usługi publikują zdarzenia domenowe. Niemutowalne fakty
opisujące to, co zaszło w~ich obszarze odpowiedzialności (np. \emph{utworzono zamówienie}). Subskrybeńci 
pobierają te zdarzenia i obsługują je według własnej logiki. Dzięki temu rozszczepieniu serwisów, nadawca nie zna swoich
odbiorców. Powoduje to łatwe dodawanie kolejnych subskrybentów bez zmian po stronie producenta \cite{kleppmann2017}.

Model ten eliminuje sprzężenie czasowe. Chwilowa niedostępność konsumenta nie wpływa na producenta.
Wiadomości nieprzetworzone będą czekały w systemie kolejkowym do ponownego działania subskrybenta. 
Ceną jest rezygnacja z~natychmiastowej, silnej spójności danych na rzecz spójności ostatecznej (\emph{eventual consistency}).
Stan rozproszony pomiędzy usługami mże być obsłużony z opóźnieniem, co system musi uwzględniać w finalnej architekturze.
Istotne są również gwarancje dostarczania. Brokery komunikatów typowo zapewniają dostarczenie \emph{co najmniej raz}, co oznacza, że ten
sam komunikat może zostać doręczony wielokrotnie. Konsumenci muszą zatem przetwarzać komunikaty w~sposób idempotentny \cite{kleppmann2017}.

W~zrealizowanym systemie oba modele zastosowano komplementarnie: żądania użytkownika obsługiwane są synchronicznie
przez REST. Natomiast proces walidacji zamówienia i~rezerwacji stanów magazynowych przebiega asynchronicznie, poprzez
zdarzenia publikowane do brokera Apache Kafka. Szczegóły tego podziału omówiono w~rozdziale~\ref{chap:projekt-systemu}.

\section{Wzorce projektowe w systemach mikroserwisowych}
\label{sec:wzorce-projektowe}

Problemy powtarzające się w~systemach mikroserwisowych doczekały się skatalogowanych rozwiązań (wzorców
projektowych) \cite{richardson2018}. Poniżej omówiono wzorce, które znalazły bezpośrednie zastosowanie
w~zrealizowanym systemie; odwołania do ich implementacji zawiera rozdział~\ref{chap:implementacja-mikroserwisow}.

\textbf{Baza danych na usługę (\emph{database-per-service}).} Każda usługa posiada wyłączność na swoje dane
i~udostępnia je jedynie poprzez interfejs. Wzorzec ten został zaimplementowany tylko częściowo gdzie 
serwis zamówień nie dotyka danych produktowych i płatniczych. Wysyłając asynchronicznie żądanie o rezerwację
produktów.

\textbf{Zdarzenia domenowe i~choreografia.} Proces biznesowy obejmujący wiele usług. Może być koordynowany
centralnie (orkiestracja) lub być realizowany jako łańcuch reakcji na zdarzenia (choreografia). Każda usługa,
przetworzywszy swój etap, publikuje zdarzenie, na które reagują kolejne. W~zrealizowanym systemie cykl walidacji
zamówienia przebiega według choreografii. Serwis zamówień publikuje zdarzenie utworzenia zamówienia, na które
serwis produktów odpowiada rezerwacją stanów magazynowych i~publikacją wyniku walidacji.

\textbf{Idempotentny odbiorca i~klucz idempotencji.} Komunikaty mogą docierać wielokrotnie, a~klient może
ponowić żądanie (np. po utracie odpowiedzi w~sieci). Operacje zmieniające stan muszą być odporne na powtórzenia.
Realizuje się to poprzez klucz idempotencji. Jest ot unikalny identyfikator operacji przekazywany przez klienta i zapamiętywany przez usługę.
Dzięki tej logice powtórne żądanie nie tworzy duplikatu, a zwraca wynik już wykonanej operacji.


\textbf{Maszyna stanów.} Cykl życia encji biznesowej w~omawianym systemie zamówienia modelowany jest jako
skończony automat stanów. Z~jawnie zdefiniowanym zbiorem dozwolonych przejść. Wzorzec ten uniemożliwia wprowadzenie
encji w~stan niespójny (np. opłacenie zamówienia anulowanego) i~porządkuje logikę biznesową wokół przejść pomiędzy
stanami.

\textbf{Brama API i~odwrócone proxy.} Klient systemu nie komunikuje się z~usługami bezpośrednio. Istnieje pojedynczy 
punkt wejścia, który przekierowuje żądania do odpowiednich serwisów na podstawie ścieżki. Upraszcza to konfiguracje 
klientów. Pozwala uniknąć CORS. Ukrywa wewnętrzną topologię systemu.

\textbf{Konfiguracja zewnętrzna.} Zgodnie z~metodyką \emph{The Twelve-Factor App} konfiguracja zależna od środowiska
(adresy usług, poświadczenia, przełączniki funkcji) przechowywana jest poza kodem. W~zmiennych środowiskowych.
Dzięki czemu ten sam artefakt może być wdrażany w~różnych środowiskach bez rekompilacji \cite{twelvefactor}.

\textbf{Kontrola zdrowia (\emph{health check}).} Każda usługa udostępnia punkt końcowy raportujący jej stan.
Odpytywany jest przez platformę uruchomieniową w~celu wykrywania niezdrowych instancji. Zapobiega to awarii
przez automatyczny ich restart lub wyłączenie z~puli.

\section{Konteneryzacja i rola orkiestracji}
\label{sec:konteneryzacja-i-orkiestracja}

Praktyczna realizacja architektury mikroserwisowej wymaga jednolitego sposobu pakowania i~uruchamiania wielu
niezależnych usług. Rolę tę pełni konteneryzacja. Jest to technika wirtualizacji na poziomie systemu operacyjnego.
Procesy izolowane są od siebie przy użyciu jądra Linux, współdzieląc przy tym jądro systemu gospodarza.
W~odróżnieniu od maszyn wirtualnych kontenery nie wymagają odrębnego systemu operacyjnego, dzięki czemu 
uruchamiają się w~ułamkach sekund i~wprowadzają znikomy narzut wydajnościowy \cite{burns2022}.

Jednostką dystrybucji jest obraz kontenera. Niemutowalny, warstwowy artefakt zawierający aplikację wraz z~kompletem
jej zależności, budowany na podstawie deklaratywnego opisu (\emph{Dockerfile}). Przchowywany jest on w~rejestrze
obrazów. Niemutowalność obrazu zapewnia powtarzalność wdrożeń. Eliminuje to błędy wynikające z~różnic z~środowiskami.

Konteneryzacja rozwiązuje problem pojedynczej usługi. Jednak nie odpowiada na pytania jak: na którym węźle uruchomić kontener?, jak utrzymać zadaną liczbę jego replik mimo awarii
węzłów?, jak kierować ruch do zmieniających się dynamicznie instancji?, jak przeprowadzić aktualizację bez przerwy w działaniu?, jak skalować liczbę replik w odpowiedzi na obciążenie.
Zadania te są określane jako orkiestracja kontenerów. Wymaga to specjalnej platformy której standardem stał się Kubernetes, któremu poświęcono następny rozdział.